\documentclass[12pt,a4paper]{report}

\usepackage{xspace}
\usepackage[square]{natbib}
\usepackage[french]{babel}
\usepackage[T1]{fontenc}
\usepackage{epsfig}
\usepackage{Damas}
\usepackage{pifont}
\usepackage{amsmath}
\usepackage{ams}
\usepackage{fancyheadings}
\usepackage{color}
\usepackage{algorithm}
\usepackage{algorithmic}

\newtheorem{theorem}{Th\'eor\`eme}
\newtheorem{lemme}{Lemme}
\newtheorem{axiom}{Axiome}
\newtheorem{note}{Note}
\newtheorem{definition}{D\'efinition}

%macros
\newcommand\pb{probl�me\xspace}
\newcommand\pbs{probl�mes\xspace}
\newcommand\ordo{ordonnancement\xspace}
\newcommand\ar{apprentissage par renforcement\xspace}
\newcommand\itema{\medskip\item}

\setcounter{secnumdepth}{3}

\textheight 9in \textwidth 145mm \topmargin -0.5in \oddsidemargin 10mm \evensidemargin 10mm \baselineskip 7mm

\lhead[\fancyplain{}{\epsfig{file=ulaval.eps,height=12pt}}]
      {\fancyplain{}{Id�es en vrac}}
\chead{\fancyplain{}{}} \rhead[\fancyplain{}{Damas}]
      {\fancyplain{}{\epsfig{file=ulaval.eps,height=12pt}}}
\lfoot[\fancyplain{}{Tous les jours, tout le temps}]
      {\fancyplain{}{Camille Besse-Patin}}
\cfoot[\fancyplain{}{\thepage}]
      {\fancyplain{}{\thepage}}
\rfoot[\fancyplain{}{Camille Besse-Patin}]
      {\fancyplain{}{Tous les jours, tout le temps}}
\setlength{\headrulewidth}{0.1pt} \setlength{\footrulewidth}{0.1pt}

\setcounter{tocdepth}{2}     % Dans la table des mati�res
\setcounter{secnumdepth}{4}  % Avec un num�ro.


\title{Id�es en vrac\,: journal de mise � jour}
\author{Camille Besse-Patin}
\director{Pr. Brahim Chaib-draa}
\encadrement{Julien Laumonier  \\ St�phane Ross}
\date{Tous les jours, tout le temps}
\typdoc{Ou comment maintenir ses id�es � jour}


\pagestyle{fancyplain}

\begin{document}

% Title Page
\pagetitre

\pagenumbering{roman}
\addtocontents{toc}{\protect\thispagestyle{fancy}}
\newpage
~\\
%Page blanche derri�re le titre
\newpage

\tableofcontents \thispagestyle{fancyplain}

\newpage
~\\

\chapter{Id�es � approfondir}
\pagenumbering{arabic} \thispagestyle{fancyplain}
\addcontentsline{toc}{chapter}{Introduction}

\section{Cymatom}


\newpage
~\\
%Page Blanche pour le d�but de chapitre en page impaire

\chapter{Id�es farfelues \& saugrenues}
\thispagestyle{fancyplain} \label{Chapter1}

\section{Comment inventer de nouvelles choses}
\begin{itemize}
    \itema Prendre un nom qui claque et essayer de faire coller une m�thode avec,
    \itema Prendre deux th�mes qui ont pas rapport et en trouver un,
    \itema Essayer de s'inspirer d'autres domaines pour trouver des algos (Cuisine, Jardinage, \ldots).
\end{itemize}

\subsection{Noms qui claque}
\begin{itemize}
    \itema Catastrophic programming (Programmation en pire cas, s�mantique d�moniaque)
    \itema Mediatic programming (Th�orie du mensonge aux agents)
    \itema Metallic programming (Recuit simul�, y'en a d'Autres � trouver peut �tre ...)
    \itema Sardonic programming (minimiser la honte/risque)
    \itema Comic programming (maximiser le fun ... TOON !!!)
    \itema Political programming
\begin{itemize}
        \itema Capitalist = individually rational
        \itema Communist = fully altruistic
        \itema Anarchist = self centered but cooperative
        \itema Democratic = hierarchical / majority votes
        \itema Dictatorial = centralized
\end{itemize}
    \itema Prismatic (d�composition lumineuse, r�fraction , r�flexion ... il y'a kk chose � creuser la).
    \itema Pathetic ... non ? ;p
\end{itemize}

\section{Comment prouver un th�or�me}

\begin{description}
  \item[by example] : \\
    The author gives only the case n = 2 and suggests that it contains most
    of the ideas of the general proof.
  \item[by intimidation] : \\
    `Trivial.'
  \item[by vigorous handwaving] : \\
    Works well in a classroom or seminar setting.
  \item[by cumbersome notation] : \\
    Best done with access to at least four alphabets and special symbols.
  \item[by exhaustion] : \\
    An issue or two of a journal devoted to your proof is useful.
  \item[by omission] : \\
    `The reader may easily supply the details.' \\
    `The other 253 cases are analogous.' \\
    `...'
  \item[by obfuscation] : \\
    A long plotless sequence of true and/or meaningless syntactically
    related statements.
  \item[by wishful citation] : \\
    The author cites the negation, converse, or generalization of a theorem
    from literature to support his claims.
  \item[by funding] : \\
    How could three different government agencies be wrong?
  \item[by eminent authority] : \\
    `I saw Karp in the elevator and he said it was probably NP-complete.'
  \item[by personal communication] : \\
    `Eight-dimensional colored cycle stripping is NP-complete
    [Karp,personal communication].'
  \item[by reduction to the wrong problem] : \\
    `To see that infinite-dimensional colored cycle stripping is decidable,
    we reduce it to the halting problem.'
  \item[by reference to inaccessible literature] : \\
    The author cites a simple corollary of a theorem to be found in a
    privately circulated memoir of the Slovenian Philological Society, 1883.
  \item[by importance] : \\
    A large body of useful consequences all follow from the proposition in
    question.
  \item[by accumulated evidence] : \\
    Long and diligent search has not revealed a counterexample.
  \item[by cosmology] : \\
    The negation of the proposition is unimaginable or meaningless. Popular
    for proofs of the existence of God.
  \item[by mutual reference] : \\
    In reference A, Theorem 5 is said to follow from Theorem 3 in reference
    B, which is shown from Corollary 6.2 in reference C, which is an easy
    consequence of Theorem 5 in reference A.
  \item[by metaproof] : \\
    A method is given to construct the desired proof. The correctness of
    the method is proved by any of these techniques.
  \item[by picture] : \\
    A more convincing form of by example. Combines well with by
    omission.
  \item[by vehement assertion] : \\
    It is useful to have some kind of authority in relation to the audience.
  \item[by ghost reference] : \\
    Nothing even remotely resembling the cited theorem appears in the
    reference given.
  \item[by forward reference] : \\
    Reference is usually to a forthcoming paper of the author, which is
    often not as forthcoming as at first.
  \item[by semantic shift] : \\
    Some standard but inconvenient definitions are changed for the
    statement of the result.
  \item[by appeal to intuition] : \\
    Cloud-shaped drawings frequently help here.
\end{description}



\end{document}
